\documentclass[header={Lecture 8 Exam Study Notes}]{study-note}

\begin{document}

\StudyTitleBlock{08}{Lecture 8: Exam Study Notes}{Hilbert spaces, projection theorem, direct sum theorem, Riesz representation, and Lax--Milgram}

\vspace{0.8em}

\begin{focusbox}
This is one of the core exam lectures. It contains the geometry of Hilbert spaces, projection/direct sum, Riesz representation, reflexivity of Hilbert spaces, and the first well-posedness result for solving equations in Hilbert spaces.

\medskip

\textbf{Convention.} Throughout these notes, the inner product $(\cdot,\cdot)_H$ is linear in the first argument and conjugate-linear in the second. This is the standard physics convention and the one used in the lectures.
\end{focusbox}

\tableofcontents

\section{Hilbert Space Basics}

\begin{defbox}{Hilbert space}
A Hilbert space is a complete inner-product space. Its norm is induced by
\[
\norm{x}_H:=\sqrt{(x,x)_H}.
\]
\end{defbox}

\begin{thmbox}{Theorems 10.1--10.3}
\thmInnerProductNorm

\medskip

\thmCauchySchwarz

\medskip

\thmParallelogram
\end{thmbox}

\section{Orthogonality, Projection, Direct Sum}

\begin{thmbox}{Theorem 11.1: Projection theorem}
If $M$ is a closed subspace of a Hilbert space $H$, then every $x\in H$ has a unique nearest point $P_Mx\in M$. Moreover,
\[
x-P_Mx\in M^\perp.
\]
\end{thmbox}

\begin{prooflevelbox}{SKETCH}
Minimize $\norm{x-y}$ over $y\in M$. Completeness gives existence of a minimizer; the parallelogram law gives uniqueness. Orthogonality: if $x-P_Mx$ had a nonzero projection onto $M$, you could move closer.
\end{prooflevelbox}

\paragraph{Broader form.}
The full projection theorem is first stated for nonempty closed convex sets. The closed-subspace version is the one used most often later, because orthogonality then characterizes the minimizer.

\begin{thmbox}{Theorem 11.2: Direct sum theorem}
\thmDirectSum
\end{thmbox}

\begin{thmbox}{Theorem 11.3: Least-squares solution in $\R^n$}
\thmLeastSquares
\end{thmbox}

\section{Riesz Representation And Reflexivity}

\begin{thmbox}{Theorem 11.4: Riesz representation theorem}
\thmRieszRepresentation
\end{thmbox}

\begin{prooflevelbox}{FULL}
Know the full proof if the theorem list or oral exam requires it: restrict $\phi$ to $\ker\phi$, use orthogonality to find the Riesz vector $a$, then show $\phi(x)=(x,a)_H$.
\end{prooflevelbox}

\paragraph{Consequence.}
Every Hilbert space is reflexive.

\paragraph{Why.}
The Riesz map $H\to H^*$, $a\mapsto (x\mapsto (x,a)_H)$, is a conjugate-linear isometry for complex scalars (and a linear isometry for real scalars). Applying Riesz to $H^*$ identifies $H^{**}$ with $H$ up to the canonical conjugation, which composes with the Riesz map to give a genuine linear isometry. Hence the canonical embedding $J:H\to H^{**}$ is onto.

\section{Lax--Milgram And Solvability}

\begin{thmbox}{Theorem 12.1: Well-posedness of $Lu=f$ under coercivity}
\thmCoerciveBijective
\end{thmbox}

\begin{prooflevelbox}{SKETCH}
Injectivity follows from coercivity. Surjectivity uses the coercivity estimate to show the range is both dense and closed.
\end{prooflevelbox}

\begin{thmbox}{Theorem 12.2: Lax--Milgram lemma}
If $B:H\times H\to\K$ is sesquilinear (linear in the first argument), bounded and coercive, then for every $f\in H^*$ there exists a unique $u\in H$ such that
\[
B(v,u)=f(v) \qquad \forall v\in H.
\]
\end{thmbox}

\begin{prooflevelbox}{SKETCH}
Know the Riesz-reduction: define $A$ by $B(v,u)=(v,Au)_H$, use coercivity to apply Theorem 12.1, then solve $Au=g$ where $g$ represents $f$ via Riesz.
\end{prooflevelbox}

\paragraph{Proof idea.}
For fixed $u$, the map $v\mapsto B(v,u)$ is in $H^*$ (linear in $v$). By Riesz, there is a unique operator $A\in\mathcal L(H)$ such that
\[
B(v,u)=(v,Au)_H.
\]
Coercivity gives
\[
\operatorname{Re}(Au,u)_H=\operatorname{Re}\ol{B(u,u)}=\operatorname{Re}B(u,u)\ge \beta\norm{u}_H^2,
\]
so Theorem \textbf{12.1} implies $A$ is invertible. Then solve $Au=g$, where $g\in H$ represents $f$ by Riesz.

\paragraph{Exam use.}
This is the abstract framework behind weak solutions of elliptic boundary value problems.

\section{Weak Convergence In Hilbert Spaces}

\begin{focusbox}
In Hilbert spaces, weak convergence can be written more concretely:
\[
x_n\rightharpoonup x \quad \Longleftrightarrow \quad (x_n,y)_H\to (x,y)_H \text{ for all } y\in H.
\]
This is just Riesz representation applied to all continuous linear functionals.
\end{focusbox}

\section{What To Memorize Precisely}

\begin{focusbox}
Know precisely:
\begin{enumerate}
\item the definition of Hilbert space and orthogonal complement;
\item the projection theorem and direct sum theorem;
\item the Riesz representation theorem;
\item why Hilbert spaces are reflexive;
\item the statements of Theorems \textbf{12.1} and \textbf{12.2};
\item how these relate to weak solutions and coercive problems.
\end{enumerate}
\end{focusbox}

\end{document}
